\documentclass[11pt,a4paper]{article}

% ============================================================
% PACOTES
% ============================================================

\usepackage[utf8]{inputenc}
\usepackage[T1]{fontenc}
\usepackage[portuguese]{babel}
\usepackage{lmodern}

\usepackage{
	amsmath,
	amssymb,
	mathtools
}

\usepackage{siunitx}
\sisetup{
	output-decimal-marker={,},
	per-mode=symbol
}

\usepackage{geometry}
\geometry{
	top=1.65cm,
	bottom=1.7cm,
	left=1.7cm,
	right=1.7cm
}

\usepackage{xcolor}
\usepackage{graphicx}
\usepackage{float}
\usepackage{booktabs}
\usepackage{array}
\usepackage{tabularx}
\usepackage{multicol}
\usepackage{enumitem}

\usepackage{tikz}
\usepackage{pgfplots}

\pgfplotsset{compat=1.18}

\usetikzlibrary{
	arrows.meta,
	positioning,
	calc,
	patterns,
	decorations.pathreplacing,
	shapes.geometric
}

\usepackage[most]{tcolorbox}
\usepackage{fancyhdr}
\usepackage{lastpage}
\usepackage{titlesec}

% ============================================================
% CORES
% ============================================================

\definecolor{eufblue}{RGB}{25,88,145}
\definecolor{euflight}{RGB}{232,243,252}
\definecolor{eufdark}{RGB}{25,48,70}
\definecolor{euforange}{RGB}{230,135,45}
\definecolor{eufgreen}{RGB}{48,137,91}
\definecolor{eufred}{RGB}{190,60,55}
\definecolor{softgray}{RGB}{246,247,249}

% ============================================================
% CONFIGURAÇÃO
% ============================================================

\setlength{\parindent}{0pt}
\setlength{\parskip}{4pt}

\setlist[enumerate]{
	itemsep=5pt,
	topsep=5pt
}

\pagestyle{fancy}
\fancyhf{}

\lhead{\textbf{\textcolor{eufblue}{EUFISICA}}}
\rhead{\textbf{Física e Química A -- 11.º ano}}

\cfoot{
	\small
	Prof. José Gonçalves
	\hfill
	Página \thepage\ de \pageref{LastPage}
}

\setlength{\headheight}{14pt}

\titleformat{\section}
{\large\bfseries\color{eufblue}}
{}
{0pt}
{}

% ============================================================
% CAIXAS
% ============================================================

\newtcolorbox{objetivos}{
	colback=euflight,
	colframe=eufblue,
	boxrule=0.9pt,
	arc=2mm,
	title=\textbf{Objetivos de aprendizagem},
	fonttitle=\bfseries
}

\newtcolorbox{recorda}{
	colback=softgray,
	colframe=eufdark,
	boxrule=0.7pt,
	arc=2mm,
	title=\textbf{Recorda},
	fonttitle=\bfseries
}

\newtcolorbox{atencao}{
	colback=euforange!8,
	colframe=euforange,
	boxrule=0.8pt,
	arc=2mm,
	title=\textbf{Atenção},
	fonttitle=\bfseries
}

\newtcolorbox{desafio}{
	colback=eufgreen!6,
	colframe=eufgreen,
	boxrule=0.8pt,
	arc=2mm,
	title=\textbf{Desafio},
	fonttitle=\bfseries
}

\newcommand{\resposta}[1][0.8]{
	\vspace{#1cm}
}

% ============================================================
% DOCUMENTO
% ============================================================

\begin{document}
	
	% ============================================================
	% CABEÇALHO
	% ============================================================
	
	\begin{center}
		
		{\Huge\bfseries\textcolor{eufblue}{EUFISICA}}
		
		\vspace{1mm}
		
		{\Large\bfseries Ficha de Trabalho}
		
		\vspace{1mm}
		
		{\LARGE\bfseries Cinemática}
		
		\vspace{1mm}
		
		{\large
			Gráficos posição--tempo, velocidade--tempo e aceleração
		}
		
		\vspace{2mm}
		
		{\small Prof. José Gonçalves}
		
	\end{center}
	
	\vspace{2mm}
	
	\begin{tabularx}{\textwidth}{|X|X|}
		\hline
		\textbf{Nome:} & \textbf{N.º:} \\[0.45cm]
		\hline
		\textbf{Turma:} & \textbf{Data:} \\[0.45cm]
		\hline
	\end{tabularx}
	
	\vspace{4mm}
	
	% ============================================================
	
	\begin{objetivos}
		
		No final desta ficha deverás ser capaz de:
		
		\begin{itemize}
			
			\item interpretar gráficos posição--tempo, velocidade--tempo e
			aceleração--tempo;
			
			\item determinar a velocidade a partir do declive de um gráfico
			posição--tempo;
			
			\item determinar a aceleração a partir do declive de um gráfico
			velocidade--tempo;
			
			\item determinar deslocamentos através da área sob um gráfico
			velocidade--tempo;
			
			\item distinguir velocidade de rapidez;
			
			\item reconhecer situações de movimento uniforme e movimento uniformemente
			variado;
			
			\item interpretar o sinal das componentes escalares da velocidade e da
			aceleração;
			
			\item representar vetorialmente a velocidade e a aceleração;
			
			\item identificar situações em que um corpo aumenta ou diminui a sua
			rapidez.
			
		\end{itemize}
		
	\end{objetivos}
	
	% ============================================================
	\section{1. Conceitos fundamentais}
	% ============================================================
	
	\begin{recorda}
		
		Num movimento retilíneo, escolhendo um eixo $Ox$:
		
		\[
		v_x=\frac{\Delta x}{\Delta t}
		\]
		
		e a aceleração média é dada por
		
		\[
		a_x=\frac{\Delta v_x}{\Delta t}.
		\]
		
		Num movimento uniformemente variado:
		
		\[
		v_x=v_{0x}+a_x t
		\]
		
		e
		
		\[
		x=x_0+v_{0x}t+\frac{1}{2}a_xt^2.
		\]
		
		Num gráfico posição--tempo, o \textbf{declive da tangente} corresponde à
		velocidade instantânea.
		
		Num gráfico velocidade--tempo:
		
		\begin{itemize}
			\item o declive corresponde à aceleração;
			\item a área algébrica entre o gráfico e o eixo do tempo corresponde
			ao deslocamento.
		\end{itemize}
		
	\end{recorda}
	
	% ============================================================
	\section{2. Interpretação de um gráfico posição--tempo}
	% ============================================================
	
	O gráfico representa o movimento retilíneo de um automóvel.
	
	\begin{figure}[H]
		\centering
		
		\begin{tikzpicture}
			
			\begin{axis}[
				width=13cm,
				height=7cm,
				axis lines=left,
				xlabel={Tempo, $t$/s},
				ylabel={Posição, $x$/m},
				xmin=0,
				xmax=14,
				ymin=-25,
				ymax=65,
				xtick={0,2,4,6,8,10,12,14},
				ytick={-20,0,20,40,60},
				grid=both,
				major grid style={gray!30},
				]
				
				\addplot[
				very thick,
				eufblue
				]
				coordinates{
					(0,0)
					(4,40)
					(7,40)
					(11,-20)
					(13,0)
				};
				
				\addplot[
				only marks,
				mark=*,
				mark size=2.5pt,
				eufblue
				]
				coordinates{
					(0,0)
					(4,40)
					(7,40)
					(11,-20)
					(13,0)
				};
				
			\end{axis}
			
		\end{tikzpicture}
		
		\caption{Posição do automóvel em função do tempo.}
		
	\end{figure}
	
	\begin{enumerate}[label=\textbf{\arabic*.}]
		
		\item Indica a posição do automóvel nos instantes:
		
		\begin{enumerate}[label=\alph*)]
			\item $t=0$;
			\item $t=\SI{4}{s}$;
			\item $t=\SI{7}{s}$;
			\item $t=\SI{11}{s}$.
		\end{enumerate}
		
		\resposta[0.8]
		
		\item Entre que instantes o automóvel esteve em repouso?
		
		\resposta[0.6]
		
		\item Em que intervalos de tempo o automóvel se deslocou no sentido positivo
		do eixo $Ox$?
		
		\resposta[0.7]
		
		\item Em que intervalo de tempo se deslocou no sentido negativo do eixo
		$Ox$?
		
		\resposta[0.7]
		
		\item Determina a componente escalar da velocidade entre
		$t=0$ e $t=\SI{4}{s}$.
		
		\resposta[1.0]
		
		\item Determina a componente escalar da velocidade entre
		$t=\SI{7}{s}$ e $t=\SI{11}{s}$.
		
		\resposta[1.0]
		
		\item Explica o significado físico do sinal negativo do resultado anterior.
		
		\resposta[1.0]
		
		\item Determina a velocidade entre $t=\SI{11}{s}$ e $t=\SI{13}{s}$.
		
		\resposta[1.0]
		
		\item Em qual dos intervalos o automóvel apresentou maior rapidez?
		
		Justifica quantitativamente.
		
		\resposta[1.2]
		
		\item Determina:
		
		\begin{enumerate}[label=\alph*)]
			\item o deslocamento total;
			\item a distância total percorrida.
		\end{enumerate}
		
		\resposta[1.4]
		
		\item Explica por que razão o deslocamento e a distância percorrida não são
		iguais neste movimento.
		
		\resposta[1.0]
		
	\end{enumerate}
	
	% ============================================================
	\section{3. Declive e velocidade}
	% ============================================================
	
	Três móveis, A, B e C, deslocam-se ao longo da mesma direção.
	
	\begin{figure}[H]
		\centering
		
		\begin{tikzpicture}
			
			\begin{axis}[
				width=12cm,
				height=7cm,
				axis lines=left,
				xlabel={$t$/s},
				ylabel={$x$/m},
				xmin=0,
				xmax=8,
				ymin=-30,
				ymax=70,
				xtick={0,1,2,3,4,5,6,7,8},
				ytick={-20,0,20,40,60},
				grid=both,
				major grid style={gray!30},
				legend style={
					at={(0.03,0.97)},
					anchor=north west
				}
				]
				
				\addplot[very thick,eufblue]
				coordinates{(0,0)(6,60)};
				\addlegendentry{A}
				
				\addplot[very thick,euforange]
				coordinates{(0,20)(6,50)};
				\addlegendentry{B}
				
				\addplot[very thick,eufred]
				coordinates{(0,40)(6,-20)};
				\addlegendentry{C}
				
			\end{axis}
			
		\end{tikzpicture}
		
		\caption{Gráficos posição--tempo dos móveis A, B e C.}
		
	\end{figure}
	
	\begin{enumerate}[label=\textbf{\arabic*.},resume]
		
		\item Calcula a componente escalar da velocidade de cada móvel.
		
		\resposta[1.5]
		
		\item Qual dos móveis apresenta maior rapidez?
		
		\resposta[0.7]
		
		\item Qual dos móveis se desloca no sentido negativo do eixo $Ox$?
		
		\resposta[0.7]
		
		\item Explica a relação entre:
		
		\begin{enumerate}[label=\alph*)]
			\item o sinal do declive e o sentido do movimento;
			\item o valor absoluto do declive e a rapidez.
		\end{enumerate}
		
		\resposta[1.3]
		
	\end{enumerate}
	
	% ============================================================
	\section{4. Da posição à velocidade}
	% ============================================================
	
	A posição de um carrinho é descrita pela expressão
	
	\[
	x(t)=4+3t,
	\]
	
	em unidades SI.
	
	\begin{enumerate}[label=\textbf{\arabic*.},resume]
		
		\item Qual é a posição inicial do carrinho?
		
		\resposta[0.5]
		
		\item Qual é a componente escalar da velocidade?
		
		\resposta[0.7]
		
		\item Em que posição se encontra o carrinho no instante
		$t=\SI{8}{s}$?
		
		\resposta[0.9]
		
		\item Representa, no referencial seguinte, o respetivo gráfico posição--tempo
		para $0\leq t\leq\SI{8}{s}$.
		
		\begin{center}
			
			\begin{tikzpicture}
				
				\begin{axis}[
					width=12cm,
					height=7cm,
					axis lines=left,
					xlabel={$t$/s},
					ylabel={$x$/m},
					xmin=0,
					xmax=9,
					ymin=0,
					ymax=32,
					xtick={0,1,2,3,4,5,6,7,8},
					ytick={0,4,8,12,16,20,24,28,32},
					grid=both,
					major grid style={gray!30}
					]
					
				\end{axis}
				
			\end{tikzpicture}
			
		\end{center}
		
		\item Representa qualitativamente o gráfico velocidade--tempo correspondente.
		
		\begin{center}
			
			\begin{tikzpicture}
				
				\begin{axis}[
					width=12cm,
					height=5cm,
					axis lines=left,
					xlabel={$t$/s},
					ylabel={$v_x$/m\,s$^{-1}$},
					xmin=0,
					xmax=9,
					ymin=-1,
					ymax=5,
					xtick={0,1,2,3,4,5,6,7,8},
					ytick={-1,0,1,2,3,4,5},
					grid=both,
					major grid style={gray!30}
					]
					
				\end{axis}
				
			\end{tikzpicture}
			
		\end{center}
		
	\end{enumerate}
	
	% ============================================================
	\section{5. Gráfico velocidade--tempo}
	% ============================================================
	
	O gráfico representa a componente escalar da velocidade de um automóvel.
	
	\begin{figure}[H]
		\centering
		
		\begin{tikzpicture}
			
			\begin{axis}[
				width=13cm,
				height=7cm,
				axis lines=middle,
				xlabel={$t$/s},
				ylabel={$v_x$/m\,s$^{-1}$},
				xmin=0,
				xmax=14,
				ymin=-12,
				ymax=22,
				xtick={0,2,4,6,8,10,12,14},
				ytick={-10,-5,0,5,10,15,20},
				grid=both,
				major grid style={gray!30}
				]
				
				\addplot[
				very thick,
				eufblue
				]
				coordinates{
					(0,0)
					(4,20)
					(7,20)
					(11,0)
					(13,-10)
				};
				
				\addplot[
				only marks,
				mark=*,
				mark size=2.5pt,
				eufblue
				]
				coordinates{
					(0,0)
					(4,20)
					(7,20)
					(11,0)
					(13,-10)
				};
				
			\end{axis}
			
		\end{tikzpicture}
		
		\caption{Componente escalar da velocidade em função do tempo.}
		
	\end{figure}
	
	\begin{enumerate}[label=\textbf{\arabic*.},resume]
		
		\item Qual é a velocidade do automóvel no instante $t=\SI{4}{s}$?
		
		\resposta[0.5]
		
		\item Entre que instantes o movimento é uniforme?
		
		\resposta[0.7]
		
		\item Entre que instantes a rapidez está a aumentar?
		
		\resposta[0.8]
		
		\item Entre que instantes a rapidez está a diminuir?
		
		\resposta[0.8]
		
		\item Em que instante o automóvel inverte o sentido do movimento?
		
		\resposta[0.7]
		
		\item Determina a aceleração entre $0$ e \SI{4}{s}.
		
		\resposta[1.1]
		
		\item Determina a aceleração entre \SI{4}{s} e \SI{7}{s}.
		
		\resposta[0.9]
		
		\item Determina a aceleração entre \SI{7}{s} e \SI{11}{s}.
		
		\resposta[1.1]
		
		\item Determina a aceleração entre \SI{11}{s} e \SI{13}{s}.
		
		\resposta[1.0]
		
		\item Determina o deslocamento do automóvel entre:
		
		\begin{enumerate}[label=\alph*)]
			
			\item $0$ e \SI{4}{s};
			
			\item \SI{4}{s} e \SI{7}{s};
			
			\item \SI{7}{s} e \SI{11}{s};
			
			\item \SI{11}{s} e \SI{13}{s}.
			
		\end{enumerate}
		
		\resposta[2.0]
		
		\item Determina o deslocamento total entre $t=0$ e $t=\SI{13}{s}$.
		
		\resposta[1.0]
		
		\item Determina a distância total percorrida.
		
		\resposta[1.0]
		
	\end{enumerate}
	
	% ============================================================
	\section{6. O significado da área num gráfico $v_x(t)$}
	% ============================================================
	
	\begin{recorda}
		
		No gráfico velocidade--tempo,
		
		\[
		\Delta x=\text{área algébrica entre a linha do gráfico e o eixo do tempo}.
		\]
		
		Áreas localizadas acima do eixo do tempo correspondem a deslocamentos
		positivos.
		
		Áreas abaixo do eixo do tempo correspondem a deslocamentos negativos.
		
	\end{recorda}
	
	Considere o gráfico seguinte.
	
	\begin{center}
		
		\begin{tikzpicture}
			
			\begin{axis}[
				width=12cm,
				height=6.5cm,
				axis lines=middle,
				xlabel={$t$/s},
				ylabel={$v_x$/m\,s$^{-1}$},
				xmin=0,
				xmax=10,
				ymin=-8,
				ymax=12,
				xtick={0,2,4,6,8,10},
				ytick={-6,-4,-2,0,2,4,6,8,10},
				grid=both,
				major grid style={gray!30}
				]
				
				\addplot[
				very thick,
				euforange
				]
				coordinates{
					(0,10)
					(4,10)
					(6,0)
					(8,-6)
					(10,-6)
				};
				
			\end{axis}
			
		\end{tikzpicture}
		
	\end{center}
	
	\begin{enumerate}[label=\textbf{\arabic*.},resume]
		
		\item Calcula o deslocamento entre $0$ e \SI{4}{s}.
		
		\resposta[0.9]
		
		\item Calcula o deslocamento entre \SI{4}{s} e \SI{6}{s}.
		
		\resposta[0.9]
		
		\item Calcula o deslocamento entre \SI{6}{s} e \SI{8}{s}.
		
		\resposta[0.9]
		
		\item Calcula o deslocamento entre \SI{8}{s} e \SI{10}{s}.
		
		\resposta[0.9]
		
		\item Determina o deslocamento total.
		
		\resposta[1.0]
		
		\item Determina a distância total percorrida.
		
		\resposta[1.0]
		
		\item Explica a diferença entre os resultados das duas questões anteriores.
		
		\resposta[1.0]
		
	\end{enumerate}
	
	% ============================================================
	\section{7. Aceleração}
	% ============================================================
	
	Um automóvel aumenta a sua velocidade de
	\SI{8}{m.s^{-1}} para \SI{20}{m.s^{-1}} durante \SI{6}{s}.
	
	\begin{enumerate}[label=\textbf{\arabic*.},resume]
		
		\item Calcula a aceleração média.
		
		\resposta[1.0]
		
		\item Interpreta fisicamente o resultado obtido.
		
		\resposta[1.0]
		
		\item Admitindo aceleração constante, determina a velocidade do automóvel
		\SI{3}{s} depois do início do intervalo considerado.
		
		\resposta[1.0]
		
	\end{enumerate}
	
	% ============================================================
	\subsection*{7.1. Uma travagem}
	% ============================================================
	
	Um automóvel desloca-se a \SI{25}{m.s^{-1}} quando o condutor inicia uma
	travagem.
	
	Ao fim de \SI{5}{s}, o automóvel encontra-se em repouso.
	
	\begin{enumerate}[label=\textbf{\arabic*.},resume]
		
		\item Determina a aceleração média durante a travagem.
		
		\resposta[1.0]
		
		\item Por que razão a aceleração apresenta sinal negativo?
		
		\resposta[0.9]
		
		\item Admitindo aceleração constante, representa qualitativamente:
		
		\begin{enumerate}[label=\alph*)]
			
			\item o gráfico velocidade--tempo;
			
			\item o gráfico aceleração--tempo.
			
		\end{enumerate}
		
		\begin{center}
			
			\begin{minipage}{0.48\textwidth}
				
				\begin{tikzpicture}
					
					\begin{axis}[
						width=\linewidth,
						height=5.5cm,
						axis lines=left,
						xlabel={$t$},
						ylabel={$v_x$},
						xmin=0,
						xmax=6,
						ymin=0,
						ymax=30,
						grid=both,
						major grid style={gray!30}
						]
						
					\end{axis}
					
				\end{tikzpicture}
				
			\end{minipage}
			\hfill
			\begin{minipage}{0.48\textwidth}
				
				\begin{tikzpicture}
					
					\begin{axis}[
						width=\linewidth,
						height=5.5cm,
						axis lines=middle,
						xlabel={$t$},
						ylabel={$a_x$},
						xmin=0,
						xmax=6,
						ymin=-8,
						ymax=4,
						grid=both,
						major grid style={gray!30}
						]
						
					\end{axis}
					
				\end{tikzpicture}
				
			\end{minipage}
			
		\end{center}
		
		\item Determina a distância percorrida durante a travagem através da área do
		gráfico velocidade--tempo.
		
		\resposta[1.2]
		
	\end{enumerate}
	
	% ============================================================
	\section{8. Velocidade e aceleração como vetores}
	% ============================================================
	
	\begin{recorda}
		
		A velocidade e a aceleração são \textbf{grandezas vetoriais}.
		
		Num movimento retilíneo:
		
		\begin{itemize}
			
			\item se $\vec{v}$ e $\vec{a}$ tiverem o mesmo sentido, a rapidez aumenta;
			
			\item se $\vec{v}$ e $\vec{a}$ tiverem sentidos opostos, a rapidez diminui.
			
		\end{itemize}
		
	\end{recorda}
	
	Analisa as quatro situações seguintes.
	
	% ------------------------------------------------------------
	
	\begin{center}
		
		\begin{tikzpicture}[>=Stealth]
			
			% A
			\node[draw,rounded corners,minimum width=6.3cm,minimum height=2.5cm]
			at (0,0) {};
			
			\node at (-2.5,0.85) {\textbf{Situação A}};
			
			\fill (-1.5,0) circle (0.12);
			
			\draw[->,very thick,eufblue]
			(-1.5,0) -- (0.3,0)
			node[midway,above] {$\vec v$};
			
			\draw[->,very thick,eufred]
			(-1.5,-0.55) -- (-0.3,-0.55)
			node[midway,below] {$\vec a$};
			
			% B
			\node[draw,rounded corners,minimum width=6.3cm,minimum height=2.5cm]
			at (7,0) {};
			
			\node at (4.5,0.85) {\textbf{Situação B}};
			
			\fill (5.5,0) circle (0.12);
			
			\draw[->,very thick,eufblue]
			(5.5,0) -- (7.3,0)
			node[midway,above] {$\vec v$};
			
			\draw[->,very thick,eufred]
			(5.5,-0.55) -- (4.3,-0.55)
			node[midway,below] {$\vec a$};
			
			% C
			\node[draw,rounded corners,minimum width=6.3cm,minimum height=2.5cm]
			at (0,-3.2) {};
			
			\node at (-2.5,-2.35) {\textbf{Situação C}};
			
			\fill (0.7,-3.2) circle (0.12);
			
			\draw[->,very thick,eufblue]
			(0.7,-3.2) -- (-1.1,-3.2)
			node[midway,above] {$\vec v$};
			
			\draw[->,very thick,eufred]
			(0.7,-3.75) -- (-0.5,-3.75)
			node[midway,below] {$\vec a$};
			
			% D
			\node[draw,rounded corners,minimum width=6.3cm,minimum height=2.5cm]
			at (7,-3.2) {};
			
			\node at (4.5,-2.35) {\textbf{Situação D}};
			
			\fill (7.8,-3.2) circle (0.12);
			
			\draw[->,very thick,eufblue]
			(7.8,-3.2) -- (6.0,-3.2)
			node[midway,above] {$\vec v$};
			
			\draw[->,very thick,eufred]
			(7.8,-3.75) -- (9.0,-3.75)
			node[midway,below] {$\vec a$};
			
		\end{tikzpicture}
		
	\end{center}
	
	\begin{enumerate}[label=\textbf{\arabic*.},resume]
		
		\item Em quais situações a rapidez está a aumentar?
		
		\resposta[0.7]
		
		\item Em quais situações a rapidez está a diminuir?
		
		\resposta[0.7]
		
		\item Na situação A, admitindo o eixo $Ox$ orientado para a direita, indica o
		sinal de:
		
		\begin{enumerate}[label=\alph*)]
			\item $v_x$;
			\item $a_x$.
		\end{enumerate}
		
		\resposta[0.7]
		
		\item Repete a questão anterior para as situações B, C e D.
		
		\resposta[1.2]
		
		\item Completa a tabela.
		
		\begin{center}
			
			\renewcommand{\arraystretch}{1.45}
			
			\begin{tabular}{c|c|c|c}
				
				\toprule
				
				\textbf{Situação} &
				\textbf{Sinal de $v_x$} &
				\textbf{Sinal de $a_x$} &
				\textbf{Rapidez}
				
				\\
				
				\midrule
				
				A & & & \\[0.35cm]
				
				B & & & \\[0.35cm]
				
				C & & & \\[0.35cm]
				
				D & & & \\[0.35cm]
				
				\bottomrule
				
			\end{tabular}
			
		\end{center}
		
	\end{enumerate}
	
	% ============================================================
	\section{9. Aceleração positiva nem sempre significa aumento da rapidez}
	% ============================================================
	
	\begin{atencao}
		
		Uma aceleração positiva não significa necessariamente que o corpo esteja a
		aumentar a sua rapidez.
		
		O efeito da aceleração depende simultaneamente dos sinais de $v_x$ e
		$a_x$.
		
	\end{atencao}
	
	Considere quatro automóveis com os seguintes valores instantâneos.
	
	\begin{center}
		
		\renewcommand{\arraystretch}{1.35}
		
		\begin{tabular}{c|c|c}
			
			\toprule
			
			\textbf{Automóvel}
			&
			$v_x$/\si{m.s^{-1}}
			&
			$a_x$/\si{m.s^{-2}}
			
			\\
			
			\midrule
			
			A & $+12$ & $+3$\\
			
			B & $+12$ & $-3$\\
			
			C & $-12$ & $+3$\\
			
			D & $-12$ & $-3$\\
			
			\bottomrule
			
		\end{tabular}
		
	\end{center}
	
	\begin{enumerate}[label=\textbf{\arabic*.},resume]
		
		\item Em quais casos a rapidez está a aumentar?
		
		\resposta[0.7]
		
		\item Em quais casos a rapidez está a diminuir?
		
		\resposta[0.7]
		
		\item Indica o sentido do movimento de cada automóvel.
		
		\resposta[0.8]
		
		\item Explica por que razão o automóvel C diminui a rapidez apesar de
		$a_x>0$.
		
		\resposta[1.0]
		
	\end{enumerate}
	
	% ============================================================
	\section{10. Movimento uniformemente acelerado}
	% ============================================================
	
	Um carrinho move-se numa trajetória retilínea segundo a equação:
	
	\[
	x(t)=2+4t+t^2
	\]
	
	em unidades SI.
	
	\begin{enumerate}[label=\textbf{\arabic*.},resume]
		
		\item Comparando com
		
		\[
		x=x_0+v_{0x}t+\frac{1}{2}a_xt^2,
		\]
		
		determina:
		
		\begin{enumerate}[label=\alph*)]
			
			\item $x_0$;
			
			\item $v_{0x}$;
			
			\item $a_x$.
			
		\end{enumerate}
		
		\resposta[1.2]
		
		\item Determina a posição do carrinho no instante $t=\SI{4}{s}$.
		
		\resposta[0.9]
		
		\item Determina a velocidade no instante $t=\SI{4}{s}$.
		
		\resposta[1.0]
		
		\item Determina o deslocamento entre $t=0$ e $t=\SI{4}{s}$.
		
		\resposta[1.0]
		
		\item Escreve a equação $v_x(t)$.
		
		\resposta[0.8]
		
		\item Representa qualitativamente os três gráficos $x(t)$, $v_x(t)$ e
		$a_x(t)$.
		
		\begin{center}
			
			\begin{minipage}{0.32\textwidth}
				
				\centering
				\textbf{$x(t)$}
				
				\begin{tikzpicture}
					
					\begin{axis}[
						width=\linewidth,
						height=5cm,
						axis lines=left,
						xlabel={$t$},
						ylabel={$x$},
						xmin=0,
						xmax=5,
						ymin=0,
						ymax=45,
						grid=both
						]
						
					\end{axis}
					
				\end{tikzpicture}
				
			\end{minipage}
			\hfill
			\begin{minipage}{0.32\textwidth}
				
				\centering
				\textbf{$v_x(t)$}
				
				\begin{tikzpicture}
					
					\begin{axis}[
						width=\linewidth,
						height=5cm,
						axis lines=left,
						xlabel={$t$},
						ylabel={$v_x$},
						xmin=0,
						xmax=5,
						ymin=0,
						ymax=16,
						grid=both
						]
						
					\end{axis}
					
				\end{tikzpicture}
				
			\end{minipage}
			\hfill
			\begin{minipage}{0.32\textwidth}
				
				\centering
				\textbf{$a_x(t)$}
				
				\begin{tikzpicture}
					
					\begin{axis}[
						width=\linewidth,
						height=5cm,
						axis lines=left,
						xlabel={$t$},
						ylabel={$a_x$},
						xmin=0,
						xmax=5,
						ymin=0,
						ymax=4,
						grid=both
						]
						
					\end{axis}
					
				\end{tikzpicture}
				
			\end{minipage}
			
		\end{center}
		
	\end{enumerate}
	
	% ============================================================
	\section{11. Relacionar os três gráficos}
	% ============================================================
	
	Um corpo apresenta o seguinte gráfico velocidade--tempo.
	
	\begin{figure}[H]
		\centering
		
		\begin{tikzpicture}
			
			\begin{axis}[
				width=12cm,
				height=6cm,
				axis lines=middle,
				xlabel={$t$/s},
				ylabel={$v_x$/m\,s$^{-1}$},
				xmin=0,
				xmax=10,
				ymin=-8,
				ymax=14,
				xtick={0,2,4,6,8,10},
				ytick={-6,-4,-2,0,2,4,6,8,10,12},
				grid=both,
				major grid style={gray!30}
				]
				
				\addplot[
				very thick,
				eufblue
				]
				coordinates{
					(0,4)
					(4,12)
					(6,12)
					(10,-4)
				};
				
			\end{axis}
			
		\end{tikzpicture}
		
	\end{figure}
	
	\begin{enumerate}[label=\textbf{\arabic*.},resume]
		
		\item Determina a aceleração entre $0$ e \SI{4}{s}.
		
		\resposta[0.9]
		
		\item Qual é a aceleração entre \SI{4}{s} e \SI{6}{s}?
		
		\resposta[0.7]
		
		\item Determina a aceleração entre \SI{6}{s} e \SI{10}{s}.
		
		\resposta[0.9]
		
		\item Em que instante, entre \SI{6}{s} e \SI{10}{s}, o corpo fica
		momentaneamente em repouso?
		
		\resposta[1.0]
		
		\item Em que intervalos:
		
		\begin{enumerate}[label=\alph*)]
			
			\item aumenta a rapidez;
			
			\item diminui a rapidez;
			
			\item apresenta rapidez constante?
			
		\end{enumerate}
		
		\resposta[1.4]
		
		\item Representa o gráfico aceleração--tempo correspondente.
		
		\begin{center}
			
			\begin{tikzpicture}
				
				\begin{axis}[
					width=12cm,
					height=6cm,
					axis lines=middle,
					xlabel={$t$/s},
					ylabel={$a_x$/m\,s$^{-2}$},
					xmin=0,
					xmax=10,
					ymin=-6,
					ymax=4,
					xtick={0,2,4,6,8,10},
					grid=both,
					major grid style={gray!30}
					]
					
				\end{axis}
				
			\end{tikzpicture}
			
		\end{center}
		
	\end{enumerate}
	
	% ============================================================
	\section{12. Da tabela para os gráficos}
	% ============================================================
	
	Num laboratório foi registada a velocidade de um carrinho.
	
	\begin{center}
		
		\renewcommand{\arraystretch}{1.35}
		
		\begin{tabular}{c|ccccccc}
			
			\toprule
			
			$t$/s & 0 & 1 & 2 & 3 & 4 & 5 & 6\\
			
			\midrule
			
			$v_x$/\si{m.s^{-1}}
			&
			2 & 4 & 6 & 8 & 10 & 12 & 14
			
			\\
			
			\bottomrule
			
		\end{tabular}
		
	\end{center}
	
	\begin{enumerate}[label=\textbf{\arabic*.},resume]
		
		\item Representa o gráfico velocidade--tempo.
		
		\begin{center}
			
			\begin{tikzpicture}
				
				\begin{axis}[
					width=12cm,
					height=6cm,
					axis lines=left,
					xlabel={$t$/s},
					ylabel={$v_x$/m\,s$^{-1}$},
					xmin=0,
					xmax=7,
					ymin=0,
					ymax=16,
					xtick={0,1,2,3,4,5,6},
					ytick={0,2,4,6,8,10,12,14,16},
					grid=both,
					major grid style={gray!30}
					]
					
				\end{axis}
				
			\end{tikzpicture}
			
		\end{center}
		
		\item Determina a aceleração do carrinho.
		
		\resposta[1.0]
		
		\item Escreve a equação da velocidade em função do tempo.
		
		\resposta[0.9]
		
		\item Determina o deslocamento entre $t=0$ e $t=\SI{6}{s}$ utilizando a área
		do gráfico.
		
		\resposta[1.3]
		
		\item Admitindo que $x_0=\SI{5}{m}$, determina a posição no instante
		$t=\SI{6}{s}$.
		
		\resposta[1.0]
		
	\end{enumerate}
	
	% ============================================================
	\section{13. Situação real -- partida de um automóvel}
	% ============================================================
	
	\begin{tcolorbox}[
		colback=euflight,
		colframe=eufblue,
		title=\textbf{Movimento num semáforo}
		]
		
		Um automóvel encontra-se parado num semáforo. Quando a luz passa a verde,
		inicia o movimento numa estrada retilínea.
		
		Durante os primeiros \SI{6}{s}, a sua aceleração é aproximadamente
		constante e igual a \SI{2.5}{m.s^{-2}}.
		
		Depois mantém a velocidade atingida durante \SI{10}{s}.
		
		Finalmente, o condutor trava uniformemente durante \SI{5}{s}, até o
		automóvel parar.
		
	\end{tcolorbox}
	
	\begin{enumerate}[label=\textbf{\arabic*.},resume]
		
		\item Qual é a velocidade ao fim dos primeiros \SI{6}{s}?
		
		\resposta[1.0]
		
		\item Representa qualitativamente o gráfico velocidade--tempo de todo o
		movimento.
		
		\begin{center}
			
			\begin{tikzpicture}
				
				\begin{axis}[
					width=13cm,
					height=7cm,
					axis lines=left,
					xlabel={$t$/s},
					ylabel={$v_x$/m\,s$^{-1}$},
					xmin=0,
					xmax=23,
					ymin=0,
					ymax=18,
					xtick={0,2,4,6,8,10,12,14,16,18,20,21},
					ytick={0,3,6,9,12,15,18},
					grid=both,
					major grid style={gray!30}
					]
					
				\end{axis}
				
			\end{tikzpicture}
			
		\end{center}
		
		\item Determina a aceleração durante a travagem.
		
		\resposta[1.0]
		
		\item Determina a distância percorrida durante os primeiros \SI{6}{s}.
		
		\resposta[1.1]
		
		\item Determina a distância percorrida durante os \SI{10}{s} de movimento
		uniforme.
		
		\resposta[1.1]
		
		\item Determina a distância percorrida durante a travagem.
		
		\resposta[1.1]
		
		\item Determina a distância total percorrida.
		
		\resposta[1.0]
		
		\item Calcula a velocidade média durante todo o percurso.
		
		\resposta[1.1]
		
	\end{enumerate}
	
	% ============================================================
	\section{14. Problema de síntese}
	% ============================================================
	
	\begin{desafio}
		
		Um móvel desloca-se numa trajetória retilínea.
		
		No instante $t=0$ encontra-se na posição $x_0=\SI{10}{m}$ e apresenta
		velocidade
		
		\[
		v_{0x}=-\SI{8}{m.s^{-1}}.
		\]
		
		A aceleração é constante e igual a
		
		\[
		a_x=+\SI{2}{m.s^{-2}}.
		\]
		
	\end{desafio}
	
	\begin{enumerate}[label=\textbf{\arabic*.},resume]
		
		\item Representa os vetores velocidade e aceleração no instante inicial,
		admitindo o eixo $Ox$ orientado para a direita.
		
		\begin{center}
			
			\begin{tikzpicture}[>=Stealth]
				
				\draw[->,thick] (-5,0)--(5,0)
				node[right] {$+x$};
				
				\fill (0,0) circle (0.12);
				
				\node[below] at (0,-0.2) {móvel};
				
			\end{tikzpicture}
			
		\end{center}
		
		\item No instante inicial, a rapidez está a aumentar ou a diminuir?
		
		Justifica.
		
		\resposta[1.0]
		
		\item Escreve a equação da velocidade:
		
		\[
		v_x(t)=\underline{\hspace{6cm}}
		\]
		
		\item Em que instante o móvel fica momentaneamente em repouso?
		
		\resposta[1.0]
		
		\item Que alteração ocorre no sentido do movimento nesse instante?
		
		\resposta[0.9]
		
		\item Determina a posição nesse instante.
		
		\resposta[1.2]
		
		\item Escreve a equação da posição em função do tempo.
		
		\resposta[1.0]
		
		\item Determina a posição no instante $t=\SI{8}{s}$.
		
		\resposta[1.0]
		
		\item Determina o deslocamento entre $t=0$ e $t=\SI{8}{s}$.
		
		\resposta[1.0]
		
		\item Determina a distância total percorrida durante esses \SI{8}{s}.
		
		\resposta[1.3]
		
		\item Esboça o gráfico velocidade--tempo entre $0$ e \SI{8}{s}.
		
		\begin{center}
			
			\begin{tikzpicture}
				
				\begin{axis}[
					width=12cm,
					height=6cm,
					axis lines=middle,
					xlabel={$t$/s},
					ylabel={$v_x$/m\,s$^{-1}$},
					xmin=0,
					xmax=9,
					ymin=-10,
					ymax=10,
					xtick={0,1,2,3,4,5,6,7,8},
					ytick={-10,-8,-6,-4,-2,0,2,4,6,8,10},
					grid=both,
					major grid style={gray!30}
					]
					
				\end{axis}
				
			\end{tikzpicture}
			
		\end{center}
		
		\item Esboça o gráfico aceleração--tempo.
		
		\begin{center}
			
			\begin{tikzpicture}
				
				\begin{axis}[
					width=12cm,
					height=5.5cm,
					axis lines=middle,
					xlabel={$t$/s},
					ylabel={$a_x$/m\,s$^{-2}$},
					xmin=0,
					xmax=9,
					ymin=-3,
					ymax=4,
					xtick={0,1,2,3,4,5,6,7,8},
					ytick={-3,-2,-1,0,1,2,3,4},
					grid=both,
					major grid style={gray!30}
					]
					
				\end{axis}
				
			\end{tikzpicture}
			
		\end{center}
		
		\item Indica qualitativamente a forma do gráfico posição--tempo.
		
		É uma reta ou uma curva? Justifica.
		
		\resposta[1.1]
		
	\end{enumerate}
	
	% ============================================================
	\section{15. Questões de escolha múltipla}
	% ============================================================
	
	\begin{enumerate}[label=\textbf{\arabic*.},resume]
		
		\item Num gráfico posição--tempo, uma reta com declive negativo significa que
		o corpo:
		
		\begin{enumerate}[label=\Alph*.]
			
			\item está necessariamente a diminuir a rapidez;
			
			\item está em repouso;
			
			\item desloca-se no sentido negativo do referencial;
			
			\item apresenta aceleração negativa.
			
		\end{enumerate}
		
		\item Num gráfico velocidade--tempo, uma linha horizontal acima do eixo do
		tempo representa:
		
		\begin{enumerate}[label=\Alph*.]
			
			\item velocidade positiva constante;
			
			\item aceleração positiva constante;
			
			\item repouso;
			
			\item velocidade crescente.
			
		\end{enumerate}
		
		\item Se $v_x<0$ e $a_x>0$, então:
		
		\begin{enumerate}[label=\Alph*.]
			
			\item a rapidez aumenta sempre;
			
			\item a rapidez diminui enquanto o corpo mantiver $v_x<0$;
			
			\item o corpo está em repouso;
			
			\item a aceleração é nula.
			
		\end{enumerate}
		
		\item A área algébrica sob um gráfico velocidade--tempo corresponde:
		
		\begin{enumerate}[label=\Alph*.]
			
			\item à aceleração;
			
			\item à rapidez;
			
			\item ao deslocamento;
			
			\item à posição inicial.
			
		\end{enumerate}
		
		\item Um automóvel apresenta $v_x=+\SI{15}{m.s^{-1}}$ e
		$a_x=-\SI{3}{m.s^{-2}}$.
		
		Nesse instante:
		
		\begin{enumerate}[label=\Alph*.]
			
			\item desloca-se no sentido positivo e aumenta a rapidez;
			
			\item desloca-se no sentido negativo e aumenta a rapidez;
			
			\item desloca-se no sentido positivo e diminui a rapidez;
			
			\item desloca-se no sentido negativo e diminui a rapidez.
			
		\end{enumerate}
		
	\end{enumerate}
	
	% ============================================================
	\section{16. Verdadeiro ou falso}
	% ============================================================
	
	Classifica cada afirmação como verdadeira (V) ou falsa (F).
	Corrige as falsas.
	
	\begin{enumerate}[label=\textbf{\arabic*.},resume]
		
		\item[$\square$]
		Num gráfico posição--tempo, a inclinação permite determinar a velocidade.
		
		\resposta[0.4]
		
		\item[$\square$]
		Uma velocidade negativa significa obrigatoriamente que o corpo está a
		abrandar.
		
		\resposta[0.5]
		
		\item[$\square$]
		Uma aceleração nula implica velocidade constante.
		
		\resposta[0.5]
		
		\item[$\square$]
		Uma aceleração negativa implica sempre diminuição da rapidez.
		
		\resposta[0.5]
		
		\item[$\square$]
		Se velocidade e aceleração tiverem o mesmo sentido, a rapidez aumenta.
		
		\resposta[0.5]
		
		\item[$\square$]
		A área sob um gráfico velocidade--tempo pode assumir um valor negativo.
		
		\resposta[0.5]
		
		\item[$\square$]
		Num gráfico velocidade--tempo, uma reta inclinada corresponde a aceleração
		não nula.
		
		\resposta[0.5]
		
	\end{enumerate}
	
	% ============================================================
	\section{17. Desafio final -- identifica o movimento}
	% ============================================================
	
	São apresentados quatro gráficos velocidade--tempo.
	
	\begin{center}
		
		\begin{tikzpicture}
			
			% gráfico A
			\begin{axis}[
				at={(0cm,6cm)},
				anchor=south west,
				width=6cm,
				height=5cm,
				axis lines=middle,
				xmin=0,
				xmax=5,
				ymin=-1,
				ymax=6,
				xlabel={$t$},
				ylabel={$v_x$},
				title={A},
				grid=both
				]
				\addplot[very thick,eufblue]
				coordinates{(0,2)(5,2)};
			\end{axis}
			
			% gráfico B
			\begin{axis}[
				at={(7cm,6cm)},
				anchor=south west,
				width=6cm,
				height=5cm,
				axis lines=middle,
				xmin=0,
				xmax=5,
				ymin=-1,
				ymax=6,
				xlabel={$t$},
				ylabel={$v_x$},
				title={B},
				grid=both
				]
				\addplot[very thick,eufgreen]
				coordinates{(0,0)(5,5)};
			\end{axis}
			
			% gráfico C
			\begin{axis}[
				at={(0cm,0cm)},
				anchor=south west,
				width=6cm,
				height=5cm,
				axis lines=middle,
				xmin=0,
				xmax=5,
				ymin=-1,
				ymax=6,
				xlabel={$t$},
				ylabel={$v_x$},
				title={C},
				grid=both
				]
				\addplot[very thick,euforange]
				coordinates{(0,5)(5,0)};
			\end{axis}
			
			% gráfico D
			\begin{axis}[
				at={(7cm,0cm)},
				anchor=south west,
				width=6cm,
				height=5cm,
				axis lines=middle,
				xmin=0,
				xmax=5,
				ymin=-6,
				ymax=2,
				xlabel={$t$},
				ylabel={$v_x$},
				title={D},
				grid=both
				]
				\addplot[very thick,eufred]
				coordinates{(0,-1)(5,-5)};
			\end{axis}
			
		\end{tikzpicture}
		
	\end{center}
	
	\vspace{3mm}
	
	Associa cada gráfico a uma das descrições seguintes.
	
	\begin{itemize}
		
		\item Movimento uniforme no sentido positivo.
		
		\item Movimento no sentido positivo com aumento da rapidez.
		
		\item Movimento no sentido positivo com diminuição da rapidez.
		
		\item Movimento no sentido negativo com aumento da rapidez.
		
	\end{itemize}
	
	\resposta[1.2]
	
	Explica, para cada gráfico, os sinais de $v_x$ e $a_x$.
	
	\resposta[1.5]
	
	% ============================================================
	\section{18. Autoavaliação}
	% ============================================================
	
	\begin{center}
		
		\renewcommand{\arraystretch}{1.45}
		
		\begin{tabularx}{\textwidth}{Xccc}
			
			\toprule
			
			\textbf{Consigo...}
			&
			\textbf{Sim}
			&
			\textbf{Com ajuda}
			&
			\textbf{Ainda não}
			
			\\
			
			\midrule
			
			Interpretar gráficos posição--tempo
			&
			$\square$
			&
			$\square$
			&
			$\square$
			\\
			
			Calcular velocidades a partir de gráficos $x(t)$
			&
			$\square$
			&
			$\square$
			&
			$\square$
			\\
			
			Interpretar gráficos velocidade--tempo
			&
			$\square$
			&
			$\square$
			&
			$\square$
			\\
			
			Determinar deslocamentos através de áreas
			&
			$\square$
			&
			$\square$
			&
			$\square$
			\\
			
			Calcular acelerações
			&
			$\square$
			&
			$\square$
			&
			$\square$
			\\
			
			Relacionar os sinais de $v_x$ e $a_x$
			&
			$\square$
			&
			$\square$
			&
			$\square$
			\\
			
			Representar vetores velocidade e aceleração
			&
			$\square$
			&
			$\square$
			&
			$\square$
			\\
			
			Relacionar os gráficos $x(t)$, $v(t)$ e $a(t)$
			&
			$\square$
			&
			$\square$
			&
			$\square$
			\\
			
			\bottomrule
			
		\end{tabularx}
		
	\end{center}
	
	\vfill
	
	% ============================================================
	% RODAPÉ FINAL
	% ============================================================
	
	\begin{center}
		
		\begin{tikzpicture}
			
			\node[
			rounded corners=3mm,
			draw=eufblue,
			line width=1pt,
			fill=euflight,
			inner xsep=10mm,
			inner ysep=4mm
			]
			{
				\begin{minipage}{0.84\textwidth}
					
					\centering
					
					{\Large\bfseries\textcolor{eufblue}{EUFISICA}}
					
					\vspace{1mm}
					
					\small
					
					Posição $\bullet$
					Velocidade $\bullet$
					Aceleração $\bullet$
					Gráficos $\bullet$
					Vetores
					
					\vspace{1mm}
					
					\textit{
						Compreender um movimento é conseguir contar a mesma história
						através de equações, vetores e gráficos.
					}
					
					\vspace{1mm}
					
					\textbf{Prof. José Gonçalves}
					
				\end{minipage}
			};
			
		\end{tikzpicture}
		
	\end{center}
	
\end{document}